\documentclass[10pt,a4paper]{amsart}
\usepackage[margin=19mm]{geometry}
\usepackage{amsmath,amssymb,mathtools}
\usepackage[scr=rsfs,cal=euler]{mathalfa}
\usepackage{xfrac}
\usepackage[unicode,hidelinks,bookmarks=false]{hyperref}
\newcommand{\ie}{i.e.\ }
\newcommand{\cf}{cf.\ }
\newcommand{\resp}{respectively\ }
\newcommand{\Subsection}{\subsection}
\providecommand{\nobreakdash}{\nobreak\hbox{-}}
\setlength{\emergencystretch}{2em}
\multlinegap0pt
\usepackage{amssymb}
\usepackage{mathtools}
\usepackage[shortlabels]{enumitem}
\newtheorem{lemma}{Lemma}
\def\R{\mathbb{R}}
\newcommand{\dGAB}{d_{\mathrm{GAB}}}
\newcommand{\dKL}{d_{\mathrm{KL}}}
\newcommand{\inner}[1]{\left\langle #1\right\rangle}
\newcommand{\norm}[1]{\left\Vert #1 \right\Vert}
\title{Characterization of Generalized Alpha-Beta Divergences and Associated Entropy Measures}
\author{Subhrajyoty Roy, Supratik Basu, Abhik Ghosh, Ayanendranath Basu}
\date{Source: arXiv:2507.04637v2 (CC BY 4.0)}
\begin{document}
\maketitle

\noindent\emph{Source and licence.} Characterization of Generalized Alpha-Beta Divergences and Associated Entropy Measures. Version arXiv:2507.04637v2 (CC BY 4.0), distributed by arXiv under the Creative Commons Attribution 4.0 International licence, http://creativecommons.org/licenses/by/4.0/, as recorded in the arXiv licence metadata for this exact version and shown on the abstract page https://arxiv.org/abs/2507.04637v2. Full licence notice: \texttt{licenses/arxiv-2507.04637-CC-BY-4.0.txt}.\par\smallskip

\noindent\textit{Editorial scope.} Counted unit: the article's lemma lem:gab-sufficient (source0.tex lines 897--899). The article's complete original proof of it is delivered with it (source0.tex lines 901--948, one \texttt{proof} environment of 48 lines). No mathematics is written, repaired or re-derived: the statement and the proof are the article's own lines, in the article's own notation, and no retained line is substituted or re-flowed.  Retained uncounted as the source-local context the proof needs: lines 245--251; lines 314--319; lines 504--507; lines 851--853.  Nothing was omitted from any retained span, so the package carries no omission record.  No retained line contains a citation, so the wrapper carries no bibliography and no bibliography entry.  No retained line contains a figure environment, an image-inclusion command or drawing code, so no image is delivered and the package declares no asset dependency.  Editorial formatting only, in addition: an AMS \texttt{amsart} preamble that restates the article's own declarations and macros used by the retained lines, namely the environments \texttt{lemma} on separate counters, together with the standard packages those lines need.  The retained environments print in this document as Lemma 1 in order of appearance, whereas the article prints the same items with the numbering of its own full document.  The complete native source of the article as distributed in the arXiv e-print for this exact version is bundled unchanged as \texttt{sources/181110/source0.tex}.
\begin{equation}
    \dGAB^{(\alpha,\beta),\psi}(P, Q)
    = \dfrac{1}{\beta(\alpha+\beta)} \psi\left( \norm{p}_{\alpha+\beta}^{\alpha+\beta} \right)
    + \dfrac{1}{\alpha(\alpha+\beta)} \psi\left( \norm{q}_{\alpha+\beta}^{\alpha+\beta} \right)
    - \dfrac{1}{\alpha\beta} \psi\left( \inner{p,q}_{\alpha,\beta} \right),
    \label{eqn:defn-gab-div}
\end{equation}

\begin{multline}
    \alpha^2 \dGAB^{(\alpha,0),\psi}(P, Q)
    = \psi'(\norm{p}_\alpha^\alpha) \norm{p}_\alpha^\alpha \left( \dKL(P^{[\alpha]}, Q^{[\alpha]}) + \ln(\norm{p}_\alpha^\alpha) - \ln(\norm{q}_\alpha^\alpha) \right)\\
    - \psi(\norm{p}_\alpha^\alpha) + \psi(\norm{q}_\alpha^\alpha).
    \label{eqn:gab-to-kl-bzero}
\end{multline}

          \begin{equation}
              \dGAB^{(\alpha,\beta),\psi}(P, Q) = \dGAB^{(\beta, \alpha),\psi}(Q, P).
              \label{eqn:gab-symmetry}
          \end{equation}

\begin{lemma}\label{lem:geometric-convexity}
    Let $\Psi(x) := \psi(e^x)$ for any $x \in \R$ and $\lambda \in [0, 1]$ be a nonnegative real number. Then $\Psi$ is $\lambda$-convex if and only if for any $x, y \geq 0$, $\lambda \psi(x) + (1 - \lambda) \psi(y) \geq \psi(x^\lambda y^{1-\lambda})$.
\end{lemma}

\begin{lemma}\label{lem:gab-sufficient}
    Define $\Psi(x) := \psi(e^x)$ for all $x \in [0, \infty)$. If $\Psi$ is strictly increasing and convex, then the GAB generated by $\psi$ is nonnegative for all choices of $\alpha$ and $\beta$ such that at least one of $\alpha$ and $\beta$ is nonzero and $(\alpha + \beta) \neq 0$.
\end{lemma}

\begin{proof}
    If $P = Q$, then it is straightforward to verify that the GAB divergence is equal to $0$. Therefore, we focus on the case when $P \neq Q$. Based on the signs of $\alpha\beta,\alpha(\alpha+\beta)$ and $\beta(\alpha+\beta)$, we deal with different cases separately using different versions of H\"{o}lder's inequality.

    \noindent\textbf{Case 1 with $\alpha\beta > 0, \alpha(\alpha+\beta) > 0, \beta(\alpha+\beta) > 0$:} By applying H\"{o}lder's inequality on $p^{\alpha+\beta}$ and $q^{\alpha+\beta}$, we have
    \begin{equation}
        \int p^\alpha q^\beta \leq \left( \int p^{\alpha+\beta} \right)^{\lambda_\alpha} \left( \int q^{\alpha+\beta} \right)^{\lambda_\beta},
        \label{eqn:thm-gab-sufficient-1}
    \end{equation}
    \noindent where $\lambda_\alpha = \alpha/(\alpha+\beta)$ and $\lambda_\beta = (1 - \lambda_\alpha)$. Since $\Psi$ is convex, by applying Lemma~\ref{lem:geometric-convexity}, we have
    \begin{equation*}
        \lambda_\alpha \psi\left( \norm{p}_{\alpha+\beta}^{\alpha+\beta}\right) + \frac{\beta}{\alpha+\beta} \psi\left( \norm{q}_{\alpha+\beta}^{\alpha+\beta}\right)
        \geq  \psi\left( \norm{p}_{\alpha+\beta}^\alpha \norm{q}_{\alpha+\beta}^\beta \right)
        \geq \psi\left( \inner{p,q}_{\alpha,\beta} \right),
    \end{equation*}
    \noindent where the last inequality follows from the inequality~\eqref{eqn:thm-gab-sufficient-1} and the strictly increasing nature of $\psi$. Dividing both sides by $\alpha\beta$ now shows that the GAB divergence form as in~\eqref{eqn:defn-gab-div} is nonnegative. Note that, the equality holds if and only if $p^{\alpha+\beta} = cq^{\alpha+\beta}$ for some constant $c > 0$. Since $p$ and $q$ are densities, then they are equal almost surely.

    \noindent\textbf{Case 2 with $\alpha\beta < 0, \alpha(\alpha+\beta) > 0, \beta(\alpha+\beta) < 0$:} Note that, $(\alpha + \beta)/\alpha > 0$, $-\beta/\alpha > 0$ and they add up to $1$. Therefore, by applying H\"{o}lder's inequality on $p^{\alpha}q^\beta$ and $q^{\alpha+\beta}$, we get
    \begin{equation*}
        \left( \int p^{\alpha} q^\beta \right)^{(\alpha+\beta)/\alpha} \left( \int q^{\alpha+\beta} \right)^{-\beta/\alpha} \geq \int p^{\alpha+\beta}.
    \end{equation*}
    \noindent Using the increasingness and geometric-convexity of $\psi$, we obtain the chain of inequalities
    \begin{equation*}
        \dfrac{\alpha+\beta}{\alpha} \psi\left( \inner{p,q}_{\alpha,\beta} \right) - \dfrac{\beta}{\alpha} \psi\left( \norm{q}_{\alpha+\beta}^{\alpha+\beta} \right)
        \geq \psi\left( \inner{p, q}_{\alpha,\beta}^{(\alpha+\beta)/\alpha} \norm{q}_{\alpha+\beta}^{-\beta(\alpha+\beta)/\alpha} \right)
        \geq \psi\left( \int p^{\alpha+\beta} \right).
    \end{equation*}

    \noindent Dividing both sides by $\beta(\alpha+\beta)$ yields,
    \begin{equation*}
        \dfrac{1}{\alpha\beta} \psi\left( \inner{p,q}_{\alpha,\beta} \right) - \dfrac{\lambda_\beta}{\alpha\beta}\psi\left( \norm{q}_{\alpha+\beta}^{\alpha+\beta} \right) \leq \dfrac{\lambda_\alpha}{\alpha\beta} \psi\left( \norm{p}_{\alpha+\beta}^{\alpha+\beta} \right).
    \end{equation*}
    \noindent A rearrangement of the above shows that the GAB divergence is nonnegative. Similar to the previous case, it is easy to see that if $p = q$, then GAB divergence is equal to $0$.

    \noindent\textbf{Case 3 with $\alpha\beta < 0, \alpha(\alpha+\beta) < 0, \beta(\alpha+\beta) > 0$:} This follows from the previous case and the duality of the GAB as in~\eqref{eqn:gab-symmetry}.

    \noindent\textbf{Case 4 with $\alpha \neq 0, \beta = 0$:} Let $p$ and $q$ be two sub-densities. Let $x = \ln(\norm{p}_{\alpha}^{\alpha})$ and $y = \ln(\norm{q}_{\alpha}^{\alpha})$, then the GAB divergence given in~\eqref{eqn:gab-to-kl-bzero} reduces to
    \begin{equation*}
        \alpha^2 \dGAB^{(\alpha,0),\psi}
        = \Psi'(x) d^\ast(P^{[\alpha]}, Q^{[\alpha]}) + \left( \Psi(y) - \Psi(x) - \Psi'(x)(y-x) \right).
    \end{equation*}
    \noindent Since, $\Psi'(x) > 0$ due to the increasing nature of $\Psi$ and the KL-divergence is nonnegative, it is enough to show that $\Psi(y)-\Psi(x)-\Psi'(x)(y-x)$ is nonnegative. It now follows from the fact that $\Psi$ is convex and $\Psi'(x)$ is a subgradient of $\Psi$ at $x$. In particular, there can be three cases. If $y = x$, then the right-hand side is $0$, hence nonnegative. If $y > x$, then by exploiting the convexity and increasing property of $\Psi$, we get $\Psi'(x)\le \Psi'(y)$ and $\Psi(x)\le \Psi(y)$. So, by Lagrange's Mean Value theorem, there exists $z\in (x,y)$ such that
    \begin{equation*}
        \Psi'(x) \leq \Psi'(z) = \frac{\Psi(y)-\Psi(x)}{y-x}\le \Psi'(y).
    \end{equation*}
    \noindent On the other hand if $y < x$, then a similar deduction applies. This completes the proof for this case.

    \noindent\textbf{Case 5 with $\alpha = 0, \beta \neq 0$:} The case for $\alpha = 0, \beta \neq 0$ is straightforward due to the symmetry of GAB divergence as in Eq.~\eqref{eqn:gab-symmetry} and the previous case.
\end{proof}

\end{document}
